% ============================================================
%  Chapitre 5 — MLOps et Pipeline IA
% ============================================================
\chapter{MLOps et Pipeline d'Intelligence Artificielle}
\thispagestyle{fancy}

\section{Introduction}

Le chapitre MLOps représente le cœur de la contribution scientifique de \farmAI{} v3.0.
Il décrit la chaîne complète de production des modèles d'IA, de l'acquisition des données
jusqu'au déploiement en production avec surveillance continue.

\section{Maturité MLOps}

\cite{sculley2015hidden} ont popularisé le concept de « dette technique cachée » dans
les systèmes ML : les modèles en production dérivent, les pipelines de données se
fragilisent, les hyperparamètres sont oubliés. Les niveaux de maturité MLOps répondent
à ces problèmes :

\begin{table}[H]
\centering
\caption{Niveaux de maturité MLOps et positionnement de FARM AI}
\label{tab:mlops_maturity}
\begin{tabular}{clll}
\toprule
\textbf{Niveau} & \textbf{Description} & \textbf{Pratiques} & \textbf{FARM AI} \\
\midrule
0 & Manuel (scripts ad-hoc)  & Pas de versionnement     & \no \\
1 & Pipeline automatisé      & DVC + reproductibilité   & \yes \\
1.5 & Tracking d'expériences & MLflow + registre modèles & \yes \\
2 & CI/CD ML complet         & GitHub Actions ML + monitoring & Partiel \\
3 & MLOps avancé (AutoML)    & Feature store + A/B testing & Hors périmètre \\
\bottomrule
\end{tabular}
\end{table}

\begin{tcolorbox}[infobox, title={Positionnement MLOps de FARM AI}]
\farmAI{} v3.0 atteint le niveau \textbf{1.5} de maturité MLOps :
pipeline DVC reproductible + tracking MLflow + registre de modèles.
Le niveau 2 (CI/CD ML automatique) est identifié comme perspective d'évolution prioritaire.
\end{tcolorbox}

\section{Pipeline DVC}

DVC (\textit{Data Version Control}) \cite{dvc2023} étend Git pour versionner les données
volumineuses et définir des pipelines reproductibles.

\subsection{Structure du pipeline}

\begin{lstlisting}[style=pythonstyle, caption={dvc.yaml — Pipeline FARM AI}, label={lst:dvc_yaml}]
stages:
  build_dataset:
    cmd: python mlops/scripts/build_dataset.py
    deps:
      - data/raw/telemetry_export.csv
      - mlops/scripts/build_dataset.py
    outs:
      - data/processed/features.csv

  train:
    cmd: python mlops/scripts/train_anomaly.py
    deps:
      - data/processed/features.csv
      - mlops/scripts/train_anomaly.py
    params:
      - mlops/params.yaml:
          - n_estimators
          - contamination
          - random_state
    outs:
      - models/anomaly_detector.pkl
    metrics:
      - mlops/metrics.json:
          cache: false
\end{lstlisting}

\subsection{Script d'entraînement avec logging MLflow}

\begin{lstlisting}[style=pythonstyle, caption={train\_anomaly.py — Stage 2 DVC avec MLflow}, label={lst:train_anomaly}]
import mlflow, joblib, json
import numpy as np
from sklearn.ensemble import IsolationForest
from sklearn.metrics import silhouette_score
import pandas as pd

FEATURES = ["soil_humidity", "pressure", "flow", "temperature"]

def main():
    df = pd.read_csv("data/processed/features.csv")[FEATURES]
    X  = df.dropna().values

    with mlflow.start_run(run_name="IsolationForest_v3"):
        # Hyperparamètres
        n_est   = int(open("mlops/params.yaml").read().split("n_estimators: ")[1].split("\n")[0])
        contam  = float(open("mlops/params.yaml").read().split("contamination: ")[1].split("\n")[0])

        mlflow.log_param("n_estimators",  n_est)
        mlflow.log_param("contamination", contam)
        mlflow.log_param("features",      FEATURES)

        # Entraînement
        clf = IsolationForest(n_estimators=n_est, contamination=contam,
                              random_state=42, n_jobs=-1)
        clf.fit(X)

        # Métriques
        labels    = clf.predict(X)          # -1 anomalie, +1 inlier
        scores    = clf.score_samples(X)
        sil       = silhouette_score(X, labels) if len(set(labels)) > 1 else 0.0
        n_anomaly = (labels == -1).sum()

        mlflow.log_metric("silhouette_score",  round(sil, 4))
        mlflow.log_metric("n_anomalies",       int(n_anomaly))
        mlflow.log_metric("anomaly_rate",      round(n_anomaly / len(X), 4))

        # Sauvegarde + enregistrement MLflow
        joblib.dump(clf, "models/anomaly_detector.pkl")
        mlflow.sklearn.log_model(clf, "model",
            registered_model_name="TelemetryAnomalyDetector")

        metrics = {"silhouette_score": sil, "n_anomalies": int(n_anomaly)}
        json.dump(metrics, open("mlops/metrics.json", "w"), indent=2)
        print(f"[MLflow] Silhouette={sil:.4f}, Anomalies={n_anomaly}/{len(X)}")

if __name__ == "__main__":
    main()
\end{lstlisting}

\section{Pipeline MLOps — Vue d'ensemble}

La figure \ref{fig:mlops_pipeline} illustre le pipeline complet, de l'ingestion des
données CSV jusqu'au service d'anomalies en production.

\begin{figure}[H]
\centering
\resizebox{\textwidth}{!}{%
% ============================================================
%  Pipeline MLOps — DVC + MLflow + IsolationForest
%  Smart Farm AI v3.0
% ============================================================
\begin{tikzpicture}[
  x=1cm, y=1cm,
  stage/.style={
    rectangle, rounded corners=5pt,
    draw=#1!70!black, fill=#1!15,
    text width=2.6cm, minimum height=1.4cm,
    align=center, font=\scriptsize\bfseries, inner sep=4pt
  },
  artifact/.style={
    rectangle, rounded corners=2pt,
    draw=farmgray!60, fill=farmgray!10,
    text width=2.2cm, minimum height=0.9cm,
    align=center, font=\tiny, inner sep=3pt,
    dashed
  },
  metric/.style={
    rectangle, rounded corners=2pt,
    draw=farmgreen!60, fill=farmgreen!10,
    text width=2cm, minimum height=0.8cm,
    align=center, font=\tiny, inner sep=3pt
  },
  arrow/.style={->, thick, >=Stealth, #1},
  darrow/.style={->, dashed, >=Stealth, farmgray!60}
]

% ─── Ligne principale du pipeline ────────────────────────────

% 1. Source données CSV (IoT)
\node[stage=farmgray] (csv) at (0, 0) {
  Données IoT\\CSV\\NODE\_A/B
};

% 2. DVC Stage 1 — build_dataset
\node[stage=farmblue] (dvc1) at (3.5, 0) {
  DVC Stage 1\\build\_dataset.py\\Normalisation
};

% 3. DVC Stage 2 — train
\node[stage=farmblue] (dvc2) at (7, 0) {
  DVC Stage 2\\train\_anomaly.py\\Entraînement
};

% 4. MLflow Run
\node[stage=farmgreen] (mlflow) at (10.5, 0) {
  MLflow Run\\Logging\\Hyperparams
};

% 5. MLflow Registry
\node[stage=farmgreen] (registry) at (14, 0) {
  MLflow\\Registry\\Model Store
};

% ─── Flèches pipeline principal ───────────────────────────────
\draw[arrow=farmblue] (csv)     -- (dvc1);
\draw[arrow=farmblue] (dvc1)    -- (dvc2);
\draw[arrow=farmgreen](dvc2)    -- node[above, font=\tiny]{log\_metrics()} (mlflow);
\draw[arrow=farmgreen](mlflow)  -- node[above, font=\tiny]{register\_model()} (registry);

% ─── Artefacts DVC (au-dessus) ────────────────────────────────
\node[artifact] (art1) at (1.75, 2.0) {
  \texttt{data/raw/}\\telemetry.csv
};
\draw[darrow] (csv.north) |- (art1.south);

\node[artifact] (art2) at (5.25, 2.0) {
  \texttt{data/processed/}\\features.csv
};
\draw[darrow] (dvc1.north) |- (art2.south);

\node[artifact] (art3) at (8.75, 2.0) {
  \texttt{models/}\\anomaly\_detector.pkl
};
\draw[darrow] (dvc2.north) |- (art3.south);

\node[artifact] (art4) at (12.25, 2.0) {
  \texttt{mlruns.db}\\648 KB
};
\draw[darrow] (mlflow.north) |- (art4.south);

% ─── Métriques MLflow (en-dessous) ───────────────────────────
\node[metric] (m_params) at (7, -2.5) {
  n\_estimators=150\\contamination=0.05\\random\_state=42
};
\draw[darrow] (dvc2.south) -- (m_params.north);

\node[metric] (m_score) at (10.5, -2.5) {
  Silhouette: 0.73\\Anomalies: 5\%\\Inliers: 95\%
};
\draw[darrow] (mlflow.south) -- (m_score.north);

% ─── Étape 6 : serve_latest_model ─────────────────────────────
\node[stage=farmorange] (serve) at (14, -3.5) {
  serve\_latest\\\_model.py\\FastAPI Service
};
\draw[arrow=farmorange] (registry.south) |-
  node[right, font=\tiny, pos=0.75]{load\_model()} (serve.north);

% ─── Étape 7 : Anomaly Alert ──────────────────────────────────
\node[stage=farmred] (alert) at (10.5, -3.5) {
  Anomaly Alert\\Dashboard\\WebSocket
};
\draw[arrow=farmred] (serve.west) -- node[above, font=\tiny]{score\_samples()} (alert.east);

% ─── Retour vers données ──────────────────────────────────────
\draw[arrow=farmgray!50, dashed, bend right=45]
  (alert.south) to node[below, font=\tiny\itshape]
  {nouvelles métriques} (csv.south);

% ─── Bandes DVC ───────────────────────────────────────────────
\draw[decorate, decoration={brace, amplitude=6pt},
      farmblue, thick] (0.5, -0.8) -- (8.5, -0.8)
  node[midway, below=8pt, font=\tiny\bfseries\color{farmblue}]
  {Pipeline DVC versionné (\texttt{dvc.yaml})};

\draw[decorate, decoration={brace, amplitude=6pt},
      farmgreen, thick] (9, -0.8) -- (15.5, -0.8)
  node[midway, below=8pt, font=\tiny\bfseries\color{farmgreen}]
  {MLflow Experiment Tracking};

% ─── Label global ─────────────────────────────────────────────
\node[font=\small\bfseries\color{farmblue}, above] at (7, 3.2)
  {Pipeline MLOps — IsolationForest Anomaly Detector};

\end{tikzpicture}

}
\caption{Pipeline MLOps — DVC, MLflow et déploiement de l'IsolationForest}
\label{fig:mlops_pipeline}
\end{figure}

\textbf{Résultats obtenus :}
\begin{itemize}
  \item Silhouette Score : \textbf{0.73} (bonne séparation inliers/anomalies)
  \item Taux d'anomalies détectées : \textbf{5\%} (paramètre \code{contamination=0.05})
  \item Taille \code{mlruns.db} : \textbf{648 KB} (historique d'expériences)
  \item Modèle enregistré : \code{TelemetryAnomalyDetector} (MLflow Registry)
\end{itemize}

\section{Vision par ordinateur — Les 8 modèles YOLO}

\farmAI{} v3.0 intègre 8 modèles YOLO v11 couvrant les principales pathologies et
espèces agricoles tunisiennes. L'architecture YOLO (\textit{You Only Look Once})
\cite{redmon2016yolo} permet une détection d'objets en temps réel dans une seule
passe forward du réseau.

\begin{table}[H]
\centering
\caption{Registre des 8 modèles YOLO intégrés dans Smart Farm AI v3.0}
\label{tab:yolo_models}
\small
\begin{tabular}{llllcc}
\toprule
\textbf{Modèle} & \textbf{Espèce/Domaine} & \textbf{Classes} &
\textbf{Architecture} & \textbf{mAP50} & \textbf{Fichier} \\
\midrule
bee\_obb      & Abeilles (ruche OBB)  & 1 classe    & YOLOv11-OBB & 0.914 & bee\_obb.pt \\
livestock     & Bétail (maladies)     & 4 classes   & YOLOv11-det & 0.891 & livestock.pt \\
leaf\_disease & Maladies foliaires    & 12 classes  & YOLOv11-det & 0.879 & leaf\_disease.pt \\
olive         & Maladies olivier      & 5 classes   & YOLOv11-det & 0.863 & olive.pt \\
insects       & Insectes ravageurs    & 8 classes   & YOLOv11-det & 0.847 & insects.pt \\
lemon         & Maladies citronnier   & 4 classes   & YOLOv11-det & 0.901 & lemon.pt \\
orange        & Maladies oranger      & 4 classes   & YOLOv11-det & 0.895 & orange.pt \\
fire\_smoke   & Feu / fumée (sécurité)& 2 classes   & YOLOv11-det & 0.934 & fire\_smoke.pt \\
\midrule
\textbf{Total} & \textbf{8 modèles}  & \textbf{40 classes} & YOLOv11 &
\textbf{0.891} & \textit{moy.} \\
\bottomrule
\end{tabular}
\end{table}

\subsection{Registre des modèles (MODEL\_REGISTRY)}

\begin{lstlisting}[style=pythonstyle, caption={MODEL\_REGISTRY — cv\_service.py}, label={lst:model_registry}]
MODEL_REGISTRY = {
    "bee": {
        "path": YOLO_PATH / "bee_obb.pt",
        "task": "obb",          # Oriented Bounding Boxes
        "classes": ["beehive"],
        "expert_profile": "bee"
    },
    "leaf_disease": {
        "path": YOLO_PATH / "leaf_disease.pt",
        "task": "detect",
        "classes": [
            "Apple__Apple_scab", "Apple__Black_rot",
            "Corn__Common_rust", "Grape__Black_rot",
            "Potato__Early_blight", "Tomato__Late_blight",
            # ... 12 classes total
        ],
        "expert_profile": "leaves"
    },
    "fire_smoke": {
        "path": YOLO_PATH / "fire_smoke.pt",
        "task": "detect",
        "classes": ["fire", "smoke"],
        "expert_profile": "fire"
    }
    # ... 5 autres modèles
}
\end{lstlisting}

\subsection{Inférence YOLO}

\begin{lstlisting}[style=pythonstyle, caption={Inférence YOLO (cv\_service.py)}, label={lst:yolo_infer}]
async def scan_image(model_key: str, image_bytes: bytes) -> dict:
    config = MODEL_REGISTRY[model_key]
    model  = YOLO(config["path"])           # chargement lazy-cached

    # Décodage image
    np_arr = np.frombuffer(image_bytes, np.uint8)
    img    = cv2.imdecode(np_arr, cv2.IMREAD_COLOR)

    # Inférence (< 200ms sur CPU)
    results = model.predict(img, conf=0.35, iou=0.45, verbose=False)

    detections = []
    for r in results:
        for box in (r.obb if config["task"] == "obb" else r.boxes):
            detections.append({
                "class":      config["classes"][int(box.cls)],
                "confidence": float(box.conf),
                "bbox":       box.xyxy[0].tolist()
            })
    return {"model": model_key, "detections": detections,
            "count": len(detections)}
\end{lstlisting}

\section{Détection d'anomalies télémétriques}

\subsection{Modèle IsolationForest}

L'IsolationForest \cite{liu2008isolation} est un algorithme de détection d'anomalies
non supervisé, particulièrement adapté aux données télémétriques IoT :

\begin{itemize}
  \item \textbf{Principe :} Isoler les anomalies par des coupes aléatoires — les
        anomalies nécessitent moins de coupes que les données normales.
  \item \textbf{Features :} \code{soil\_humidity}, \code{pressure}, \code{flow},
        \code{temperature} (NODE\_A) ; \code{weight}, \code{hive\_temp},
        \code{ext\_hum} (NODE\_B)
  \item \textbf{Configuration :} \code{n\_estimators=150}, \code{contamination=0.05},
        \code{random\_state=42}
\end{itemize}

\section{Assistant Souverain — IA Générative en Darija}

\subsection{Architecture du service RAG}

L'assistant de \farmAI{} v3.0 combine trois composants :

\begin{enumerate}
  \item \textbf{LLM Labess-7B} \cite{labess7b} — modèle de langage spécialisé en
        dialecte arabe tunisien, exécuté localement via Ollama.
  \item \textbf{RAG ChromaDB} \cite{chromadb2023} — base vectorielle contenant des
        documents officiels tunisiens (UTAP, AVFA) encodés avec \code{sentence-transformers}.
  \item \textbf{LLaVA Vision} \cite{llava2023} — modèle multimodal pour l'analyse
        d'images ESP32-CAM.
\end{enumerate}

\begin{lstlisting}[style=pythonstyle, caption={Service RAG — rag\_service.py (extrait)}, label={lst:rag}]
import chromadb
from sentence_transformers import SentenceTransformer

encoder    = SentenceTransformer("sentence-transformers/paraphrase-multilingual-mpnet-base-v2")
chroma_cli = chromadb.PersistentClient(path="./chromadb_data")
collection = chroma_cli.get_or_create_collection(
    name="sovereign_wisdom_v3",
    metadata={"hnsw:space": "cosine"}
)

def query_rag(question: str, n_results: int = 5) -> list[str]:
    embedding = encoder.encode([question]).tolist()
    results   = collection.query(
        query_embeddings=embedding,
        n_results=n_results,
        include=["documents", "metadatas"]
    )
    return results["documents"][0]
\end{lstlisting}

\subsection{Profils d'experts}

L'assistant utilise 7 profils d'experts adaptés à l'espèce détectée par YOLO :

\begin{table}[H]
\centering
\caption{Profils d'experts de l'Assistant Souverain}
\label{tab:expert_profiles}
\begin{tabular}{lll}
\toprule
\textbf{Profil} & \textbf{Déclencheur} & \textbf{Spécialité} \\
\midrule
\code{bee}      & Modèle bee\_obb      & Apiculture, Varroa, pathologies abeilles \\
\code{cow}      & Modèle livestock (bovins) & Pathologies bovines \\
\code{sheep}    & Modèle livestock (ovins)  & Pratiques d'élevage ovin \\
\code{leaves}   & Modèle leaf\_disease & Maladies foliaires, fongicides \\
\code{olive}    & Modèle olive         & Maladies de l'olivier, taille \\
\code{fire}     & Modèle fire\_smoke   & Premiers secours incendie, alerte \\
\code{general}  & Défaut               & Agriculture générale, Darija \\
\bottomrule
\end{tabular}
\end{table}

\subsection{Mode Lite (Oracle Cloud sans GPU)}

Pour le déploiement sur Oracle Cloud Free Tier (ARM64, sans GPU), un mode
\code{LITE\_MODE=true} remplace l'inférence LLM par une base de connaissances
synthétique statique :

\begin{lstlisting}[style=pythonstyle, caption={Lite Mode — fallback sans LLM}, label={lst:lite_mode}]
if settings.LITE_MODE:
    # Pas d'Ollama disponible — réponse depuis KB statique
    response = STATIC_KB.get(profile, {}).get(
        "response",
        "Veuillez consulter un vétérinaire pour ce cas."
    )
else:
    # Ollama local → Groq API → static fallback
    context  = query_rag(question)
    prompt   = build_prompt(profile, question, context)
    response = ollama_generate(prompt)
\end{lstlisting}
